% Chapter 5: Analysis — ablation, cold-start, hard-residual, efficiency, error analysis
\selectlanguage{greek}


\chapter{Ανάλυση}
\label{ch:analysis}

\section{Στρέβλωση στη Δειγματοληψία Αρνητικών Παραδειγμάτων}
\label{sec:hub-artifact}

Η βαθμολογία \en{Macro F1} = \CascadeValFone{} του \en{CascadeLP}
(\S\ref{sec:cascade-results}) δεν εξηγείται από απλή επικάλυψη ζευγών
εκπαίδευσης και ελέγχου (\S\ref{sec:leakage-results}). Αυτός ο έλεγχος όμως
δεν καλύπτει στρεβλώσεις στη διαδικασία δειγματοληψίας ούτε τη διαρροή
δομικής πληροφορίας από την κατασκευή του γραφήματος πριν από τη διαίρεση
(\S\ref{subsec:data-splits}). Μια βαθμολογία τόσο κοντά στο τέλειο σε ένα πρόβλημα διμελούς
ταξινόμησης 948 χιλιάδων ζευγών αξίζει έναν επιπλέον, ανεξάρτητο έλεγχο: πόσο
από αυτή την ακρίβεια θα μπορούσε να αναπαραχθεί χωρίς καμία σχεσιακή ή
σημασιολογική πληροφορία;

Ομαδοποιώντας τις μη αυτοαναφορικές γραμμές του \texttt{\en{train.csv}} ανά
τιμή του πρώτου άκρου \texttt{\en{id1}}, βρίσκουμε ότι από \HubUniqueSourceCount{}
μοναδικές τιμές, ακριβώς \HubCount{} εμφανίζονται τουλάχιστον \HubThreshold{}
φορές η καθεμία---το ίδιο σύνολο τιμών προκύπτει και σε όρια 50 και 200, οπότε
το εύρημα δεν εξαρτάται από την ακριβή επιλογή ορίου.
Αυτοί οι \HubCount{}
«κόμβοι-συγκέντρωσης» αντιστοιχούν σε \HubRowCount{} γραμμές (\HubRowPct\% του
συνόλου) και είναι \HubPurityPct\% αρνητικοί.
Κάθε άλλη γραμμή---\NonHubRowCount{}
γραμμές---είναι θετική με ακρίβεια \NonHubPosPct\%.
Ένας ταξινομητής
αναζήτησης---πρόβλεψη της πλειοψηφικής ετικέτας του κόμβου-συγκέντρωσης όταν
το \texttt{\en{id1}} ανήκει σε αυτούς τους \HubCount{}, αλλιώς πρόβλεψη θετικής
ετικέτας---χωρίς κείμενο, γράφημα, ή εκπαίδευση μοντέλου, πετυχαίνει
\en{Macro F1} = \HubLookupFone{} και ακρίβεια \HubLookupAcc{} σε ολόκληρο το
\texttt{\en{train.csv}}.
Και οι \HubTestOverlap{} τιμές εμφανίζονται στο
\texttt{\en{test.csv}}, καλύπτοντας το \HubTestRowPct\% των γραμμών ελέγχου---
η ίδια συντόμευση είναι επομένως διαθέσιμη και στο πραγματικό σύνολο ελέγχου
και στον δημόσιο πίνακα κατάταξης του διαγωνισμού.

\begin{figure}[htbp]
  \centering
  \includegraphics[width=\textwidth]{negative_sampling_artifact}
  \caption{Η στρέβλωση στη δειγματοληψία αρνητικών παραδειγμάτων στο \en{DSAA 2023}. Οι
  \HubCount{} συχνές τιμές του πρώτου άκρου είναι σχεδόν αποκλειστικά
  αρνητικές, ενώ όλες οι υπόλοιπες είναι θετικές, και το ίδιο σύνολο πηγών
  καλύπτει παρόμοιο ποσοστό του αρχείου ελέγχου. Ο κανόνας αναζήτησης μόνο από
  την ταυτότητα του πρώτου άκρου πετυχαίνει \en{Macro F1} =
  \HubLookupFone{}.}
  \label{fig:negative-sampling-artifact}
\end{figure}

Ένα μόνο παράδειγμα καθιστά το φαινόμενο απτό.
Το άρθρο \en{«2016 Chinese FA
Women's Super Cup»} (ένα ποδοσφαιρικό κύπελλο) είναι κόμβος-συγκέντρωσης,
συζευγμένο στο σύνολο εκπαίδευσης με εκατοντάδες ουσιαστικά τυχαίους στόχους,
όλους με αρνητική ετικέτα: το έντομο \en{«E. placida»}, τον ηθοποιό \en{«Frank
Jude Jr.»}, τον αστρονόμο του 17ου αιώνα \en{«Francesco Stelluti»}, το
αεροδρόμιο \en{«Omaka Aerodrome»}, και άλλους.
Καμία θεματική, γεωγραφική, ή
χρονική σχέση δεν συνδέει κανένα από αυτά τα ζεύγη με το άρθρο του κυπέλλου---
πρόκειται ορατά για έναν κόμβο-πηγή συζευγμένο με ένα αυθαίρετο δείγμα του
σώματος κειμένων, όχι για μια επιμελημένη συλλογή εύλογων μη-συνδέσμων.
Ο
πληθυσμός εκτός κόμβων-συγκέντρωσης έχει διάμεσο πλήθος μίας γραμμής ανά τιμή
\texttt{\en{id1}} και σχεδόν καθολική θετική ετικέτα---συνεπές με δείγμα
γνήσιων υπερσυνδέσμων.
Ο πληθυσμός των κόμβων-συγκέντρωσης έχει έως και
$\sim$8.000 γραμμές ανά τιμή---συνεπές με μαζική δειγματοληψία αρνητικών
παραδειγμάτων γύρω από έναν μικρό αριθμό αυθαίρετων άρθρων-αγκυρών.

\subsection{Το Τεχνούργημα «εν Δράσει»}
\label{subsec:hub-in-action}

Το παραπάνω εύρημα αφορά το σύνολο δεδομένων καθαυτό.
Για να δείξουμε ότι
επηρεάζει και το πραγματικό, ήδη αναφερόμενο μοντέλο, διασταυρώνουμε το
σύνολο των \HubCount{} κόμβων-συγκέντρωσης απευθείας με τις υπαρκτές
προβλέψεις του \en{CascadeLP} στο σύνολο επικύρωσης
(\texttt{\en{outputs\slash{}predictions\slash{}cascade\_val\_tiers.csv}},
χωρίς καμία επανεκπαίδευση ή αλλαγή πρωτοκόλλου).
Το \HubValSourcedPct\% των
\ValSplitSize{} ζευγών επικύρωσης έχει κόμβο-συγκέντρωσης ως πρώτο άκρο· ένας
απλός πίνακας αναζήτησης πάνω σε αυτό ακριβώς το υποσύνολο ζευγών πετυχαίνει
\en{Macro F1} = \HubValLookupFone{}.
Ο κανόνας ορίζεται περιγραφικά από το πλήρες επισημασμένο αρχείο και επομένως
δεν αποτελεί ανεξάρτητα εκπαιδευμένο μοντέλο αναφοράς.
Χρησιμοποιείται ως
διαγνωστικός έλεγχος του πόσο άμεσα κωδικοποιείται η ετικέτα στην ταυτότητα
του πρώτου άκρου.
Πιο αποκαλυπτικό: οι προβλέψεις του πραγματικού, εκπαιδευμένου
\en{CascadeLP} συμφωνούν με τον τετριμμένο πίνακα αναζήτησης στο
\HubValAgreementPct\% όλων των ζευγών επικύρωσης, και στο
\HubValHubAgreementPct\% των ζευγών με κόμβο-συγκέντρωσης.
Το εκπαιδευμένο
μοντέλο, δηλαδή, δεν διαφέρει σχεδόν καθόλου από τον απλό κανόνα «είναι το
\texttt{\en{id1}} ένας από αυτούς τους \HubCount{} κόμβους;» σε πάνω από τα
μισά ζεύγη του συνόλου επικύρωσης.

Αυτό δεν σημαίνει ότι το \en{CascadeLP} είναι λανθασμένο ως αρχιτεκτονική---
δρομολογεί σωστά και κλιμακώνει με βάση την αξιοπιστία των υποκείμενων
χαρακτηριστικών, όπως σχεδιάστηκε.
Σημαίνει ότι το συγκεκριμένο σύνολο
δεδομένων \en{DSAA 2023} δεν επιτρέπει να διακρίνουμε «το \en{CascadeLP}
έμαθε να προβλέπει υπερσυνδέσμους» από «το \en{CascadeLP} έμαθε να αναγνωρίζει
\HubCount{} συγκεκριμένους κόμβους»---και οι δύο υποθέσεις προβλέπουν σχεδόν
πανομοιότυπα αποτελέσματα σε αυτό το σύνολο δεδομένων.
Η ίδια ασάφεια πλήττει
πιθανότατα και κάθε άλλη αναφερόμενη σχεδόν τέλεια βαθμολογία σε αυτό το
σύνολο δεδομένων, συμπεριλαμβανομένων προηγούμενων υποβολών του διαγωνισμού
που βασίζονται σε χαρακτηριστικά \en{POS} (Κεφάλαιο~\ref{ch:background}): καμία
δεν αναφέρει έλεγχο ανάλογο του παρόντος.
Η \S\ref{sec:wiki-fresh} απαντά στο
ερώτημα που αφήνει ανοιχτό αυτή η ασάφεια, εκπαιδεύοντας από την αρχή την ίδια
αρχιτεκτονική \en{CascadeLP} σε ένα σύνολο δεδομένων του οποίου η διαδικασία τυχαίας
δειγματοληψίας αρνητικών αποφεύγει τη συγκεκριμένη συγκέντρωση πηγών.

Η υψηλή απόδοση του \texttt{\en{PosClassifier}} στο \en{DSAA 2023}
(\en{Macro F1} = \TierTwoFone{} ως αυτόνομο μοντέλο αναφοράς,
\S\ref{sec:cascade-results}, χειριζόμενο το \TierTwoPct\% του συνόλου
επικύρωσης) χρειάζεται την ίδια επιφύλαξη με τη συνολική βαθμολογία του
\en{CascadeLP}.
Ένα \en{Random Forest} πάνω σε συχνότητες μερών του λόγου δεν
έχει προφανή λόγο να αναγνωρίζει την ύπαρξη υπερσυνδέσμου με τέτοια ακρίβεια·
όμως αν οι κόμβοι-συγκέντρωσης του \texttt{\en{id1}} έχουν συστηματικά
διαφορετικό ύφος κειμένου (π.χ. άρθρα αθλητικών διοργανώσεων, όπως στο
παράδειγμα \en{«2016 Chinese FA Women's Super Cup»} παραπάνω) από τον
πληθυσμό-στόχο με τον οποίο συζευγνύονται, το διάνυσμα \en{POS} μπορεί να
«αναγνωρίζει» τον χαρακτηριστικό κόμβο-πηγή, όχι τη σχέση μεταξύ των δύο
άκρων --- η ίδια στρέβλωση σε διαφορετική οικογένεια
χαρακτηριστικών.
Η σχεδόν τέλεια βαθμολογία δεν αποδεικνύει άρα ότι το μοντέλο
έμαθε μια γνήσια σχέση μεταξύ δύο άρθρων· είναι ακριβώς αυτό που θα αναμενόταν
αν το μοντέλο εκμεταλλεύεται τη στρέβλωση στη δειγματοληψία.
Η αντίστιξη με το
ανεξάρτητο σύνολο Απλής Αγγλικής \en{Wikipedia} είναι εδώ κατατοπιστική: υπό
την ίδια αρχιτεκτονική και τις ίδιες προεπιλεγμένες υπερπαράμετρους, ο
αυτόνομος \texttt{\en{PosClassifier}} πέφτει σε \en{Macro F1} = \WfPosFone{}
(\S\ref{sec:wiki-fresh}), και η κατάταξη των οικογενειών χαρακτηριστικών
αλλάζει ουσιωδώς --- συνεπές με την υπόθεση ότι η επίδοσή του στο \en{DSAA
2023} αντανακλά σε μεγάλο βαθμό αυτή τη στρέβλωση και όχι εγγενή ισχύ της
οικογένειας χαρακτηριστικών \en{POS} καθαυτής.


\section{Ευαισθησία στη Σειρά των Άκρων}
\label{sec:endpoint-swap}

Η ετικέτα-στόχος αυτής της εργασίας είναι συμμετρική ως προς τα δύο άκρα ενός
ζεύγους: $y_{uv} = y_{vu}$ εξ ορισμού (\S\ref{sec:problem-definition}). Ο
\texttt{\en{PosClassifier}} του Επιπέδου~2 όμως συνενώνει τα δύο 36-διάστατα
ιστογράμματα \en{POS} διατηρώντας τη σειρά $[\mathbf{p}_u \,\|\, \mathbf{p}_v]$
αντί να τα συμμετροποιεί (\S\ref{subsec:pos-features}), οπότε τίποτα στην
αρχιτεκτονική δεν εγγυάται ότι η πρόβλεψη παραμένει αναλλοίωτη αν ανταλλαχθούν
τα $u$ και $v$.
Ποσοτικοποιούμε την ευαισθησία αυτή αντιστρέφοντας κάθε ζεύγος
του συνόλου επικύρωσης --- εναλλάσσοντας τα δύο μπλοκ χαρακτηριστικών \en{POS}
και τα αντίστοιχα πεδία \texttt{\en{id1}}/\texttt{\en{id2}}, διατηρώντας
αμετάβλητα τα συμμετρικά δομικά και \en{Sentence-Transformer} χαρακτηριστικά
--- και τροφοδοτώντας το ζεύγος στο ήδη εκπαιδευμένο, αναλλοίωτο
\texttt{\en{cascade.joblib}}, χωρίς καμία επανεκπαίδευση
(\texttt{\en{scripts\slash{}dsaa\slash{}audit\_protocol.py --swap}}). Πριν από
τη σύγκριση, ο έλεγχος επαληθεύει ότι οι προβλέψεις και οι δρομολογήσεις πάνω
στα μη-αντεστραμμένα ζεύγη αναπαράγουν επακριβώς το αποθηκευμένο
\texttt{\en{cascade\_val\_tiers.csv}} (\S\ref{subsec:implementation-architecture}),
ώστε η σύγκριση να αφορά πράγματι το ίδιο μοντέλο και το ίδιο σύνολο.

Σε \SwapN{} ζεύγη επικύρωσης, η αντιστροφή αλλάζει την τελική πρόβλεψη του
\en{CascadeLP} στο \SwapLabelChangePct\% των περιπτώσεων, με μέση απόλυτη
μεταβολή πιθανότητας \SwapMeanAbsProbChange{} και αλλαγή του επιπέδου που
αποφασίζει σε \SwapTierChanges{} ζεύγη.
Η επίδραση είναι έντονα ασύμμετρη ως
προς την αρχική ετικέτα: στα ζεύγη με αρχική αρνητική ετικέτα η αντιστροφή
αλλάζει την πρόβλεψη στο \SwapNegClassChangePct\% των περιπτώσεων, ενώ στα
ζεύγη με αρχική θετική ετικέτα μόλις στο \SwapPosClassChangePct\%.
Με άλλα
λόγια, οι προβλέψεις πάνω σε (κατά κανόνα δειγματοληπτημένα) αρνητικά ζεύγη
είναι σχεδόν πάντα ασταθείς ως προς τη σειρά των άκρων, ενώ οι προβλέψεις πάνω
σε γνήσια θετικά ζεύγη παραμένουν σε μεγάλο βαθμό σταθερές.

\begin{figure}[htbp]
  \centering
  \includegraphics[width=0.8\textwidth]{endpoint_swap_sensitivity}
  \caption{Ποσοστό αλλαγής πρόβλεψης μετά την αντιστροφή των άκρων. Η μεγάλη
  διαφορά μεταξύ αρχικά αρνητικών και θετικών ζευγών συνδέει την αστάθεια του
  διατεταγμένου διανύσματος \en{POS} με τη θέση του πρώτου άκρου στο
  \en{DSAA 2023}.}
  \label{fig:endpoint-swap}
\end{figure}

Το εύρημα αυτό \textbf{δεν} τεκμηριώνει ότι το πρόβλημα θα έπρεπε να
αντιμετωπιστεί ως κατευθυνόμενο --- η ετικέτα παραμένει συμμετρική εξ ορισμού
(\S\ref{sec:problem-definition}), και η εργασία δεν αλλάζει τον ορισμό του
προβλήματος με βάση αυτό το εύρημα.
Τεκμηριώνει αντίθετα ένα συγκεκριμένο
περιορισμό \textit{υλοποίησης και αξιολόγησης}: ο \texttt{\en{PosClassifier}}
μπορεί να εκμεταλλεύεται την αυθαίρετη θέση ενός άκρου (ποιο από τα δύο
καταγράφεται ως \texttt{\en{id1}}) αντί μόνο τη σχέση μεταξύ των δύο άρθρων.
Η έντονη ασυμμετρία ανά κλάση συνδέεται εύλογα με τη στρέβλωση της
\S\ref{sec:hub-artifact}: αφού οι κόμβοι-συγκέντρωσης εμφανίζονται σχεδόν
αποκλειστικά ως πρώτο άκρο σε αρνητικά ζεύγη, η αντιστροφή τους στη θέση του
δεύτερου άκρου αφαιρεί ακριβώς το σήμα θέσης που ο ταξινομητής είχε μάθει να
συσχετίζει με αρνητική ετικέτα --- συνεπές αλλά ανεξάρτητο εύρημα από τον
έλεγχο της \S\ref{subsec:hub-in-action}, δεδομένου ότι η αντιστροφή δεν
αλλάζει καθόλου ποια τιμή \texttt{\en{id1}}/\texttt{\en{id2}} ανήκει στο σύνολο
των κόμβων-συγκέντρωσης, μόνο σε ποια θέση του διανύσματος \en{POS}
εμφανίζεται.
Δεν αποκλείεται η ασυμμετρία να οφείλεται εν μέρει και σε άλλες
συστηματικές διαφορές μεταξύ πρώτου και δεύτερου άκρου πέραν της ταυτότητας
κόμβου-συγκέντρωσης· ο παρών έλεγχος δεν τις απομονώνει.
Καμία επανεκτέλεση με συμμετροποιημένο διάνυσμα \en{POS} (π.χ. άθροισμα ή
απόλυτη διαφορά των δύο ιστογραμμάτων αντί για συνένωση με σειρά) δεν
εκτελέστηκε στην παρούσα φάση· παραμένει άμεση κατεύθυνση μελλοντικής εργασίας
(Κεφάλαιο~\ref{ch:conclusion}).


\section{Ανάλυση Ευαισθησίας των Ορίων Δρομολόγησης}
\label{sec:ablation}

Σαρώνουμε ένα πλέγμα τιμών $\tau_1 \in \{0{,}6, 0{,}7, 0{,}8, 0{,}9, 0{,}95\}$ και
$\tau_2 \in \{0{,}5, 0{,}6, 0{,}7, 0{,}8, 0{,}9\}$ (\AblGridPoints{} συνδυασμοί) στο ήδη εκπαιδευμένο
\en{CascadeLP}, χωρίς επανεκπαίδευση των ταξινομητών ανά επίπεδο --- τα όρια
επηρεάζουν μόνο τη δρομολόγηση κατά το συμπερασμό (\S\ref{sec:cascade-routing}) --- και
καταγράφουμε το \en{Macro F1} και το ποσοστό δρομολόγησης ανά επίπεδο στο σύνολο
επικύρωσης για κάθε συνδυασμό.

Το όριο $\tau_1$ παρουσιάζει μια απότομη ασυνέχεια: για $\tau_1 \in \{0{,}6, 0{,}7\}$
το Επίπεδο~1 χειρίζεται \AblTierOneLowPct\% των ζευγών (\en{Macro F1} \AblTierOneLowFone{}--\AblTierOneHighFone), ενώ
για $\tau_1 \geq 0{,}8$ το ποσοστό αυτό καταρρέει απότομα στο \AblTierOneBreakPct\% (και περαιτέρω
στο \AblTierOneBreakPctHigh\% για $\tau_1 \geq 0{,}9$), με ταυτόχρονη \textit{βελτίωση} του συνολικού
\en{Macro F1} στο \AblHighTauOneLowFone{}--\AblBestFone.
Αυτό είναι αναμενόμενο δεδομένης της ακραίας
ανισορροπίας κλάσεων στο υποσύνολο προβλέψεων υψηλής βεβαιότητας του Επιπέδου~1
(\S\ref{sec:cascade-results}): χαλαρώνοντας το όριο επιτρέπουμε σε περισσότερα ζεύγη
να επιλυθούν από έναν ταξινομητή που ουσιαστικά προβλέπει πάντα τη θετική ετικέτα σε
αυτό το καθεστώς, αυξάνοντας τον κίνδυνο σφάλματος στην σπάνια αρνητική κλάση.

Το όριο $\tau_2$ έχει μονότονη επίδραση: αυξάνοντάς το από 0{,}5 σε 0{,}9 (σταθερού
$\tau_1$), το \en{Macro F1} μειώνεται σταδιακά --- π.χ. από \AblBestFone{} σε \AblHighTauOneLowFone{} για
$\tau_1 = 0{,}8$ --- ενώ το ποσοστό δρομολόγησης στο Επίπεδο~3 αυξάνεται από 0\% σε \AblTauTwoMaxTierThreePct\%.
Αυτό αντικατοπτρίζει άμεσα την αναξιοπιστία του \texttt{\en{EmbeddingClassifier}} στο
υπόλειμμα ζευγών που φτάνει ως εκεί (\en{Macro F1} \TierThreeFone, \S\ref{sec:hard-residual}):
κάθε επιπλέον ζεύγος που κλιμακώνεται στο Επίπεδο~3 αντί να επιλυθεί στο Επίπεδο~2
είναι, κατά μέσο όρο, ζημιογόνο για τη συνολική βαθμολογία.
Ο καλύτερος συνδυασμός στο
πλέγμα είναι $\tau_1 \in \{0{,}8, 0{,}9, 0{,}95\}$ με $\tau_2 = 0{,}5$
(\en{Macro F1} = \AblBestFone{}, μηδενικό ποσοστό δρομολόγησης στο Επίπεδο~3), ελαφρώς πάνω από τους
προεπιλεγμένους $\tau_1 = \TauOneDefault$, $\tau_2 = \TauTwoDefault$ (\en{Macro F1} = \AblDefaultFone{}). Η
διαφορά είναι μικρή (0{,}0001), αλλά δεν υπολογίστηκε διάστημα αβεβαιότητας ή
επανάληψη με διαφορετικές διαιρέσεις ώστε να χαρακτηριστεί ως θόρυβος.
Διατηρούμε
την προεπιλεγμένη διαμόρφωση και αντιμετωπίζουμε τη σάρωση ως περιγραφική ανάλυση
ευαισθησίας· η εμφανής εξαίρεση είναι η απότομη ασυνέχεια στο
$\tau_1 \approx 0{,}75$.
Το πλέγμα αυτό περιορίζεται σκόπιμα σε $\tau_2 \leq 0{,}9$· η \S\ref{subsec:force-tier3} επεκτείνει
τη σάρωση πολύ πέρα από αυτό το εύρος για να εξετάσει ρητά το σενάριο όπου το Επίπεδο~3
εξαναγκάζεται να χειριστεί πολύ μεγαλύτερο μερίδιο ζευγών.

\subsection[Εξαναγκασμός Περισσότερων Ζευγών στο Επίπεδο 3]{Εξαναγκασμός Περισσότερων Ζευγών στο Επίπεδο~3}
\label{subsec:force-tier3}

Το πλέγμα της \S\ref{sec:ablation} περιορίζεται σε $\tau_2 \leq 0{,}9$ επειδή αυτό είναι
το ρεαλιστικό εύρος λειτουργίας --- πέρα από αυτό, ρωτάμε ένα διαφορετικό ερώτημα: πόσο
μακριά μπορεί να φτάσει το ποσοστό κλήσεων του Επιπέδου~3 αν ο μοναδικός στόχος είναι να
δρομολογηθούν όσο το δυνατόν περισσότερα ζεύγη σε αυτό, και τι κερδίζεται ή χάνεται όταν
συμβεί αυτό; Επεκτείνουμε τη σάρωση του $\tau_2$ στο $\tau_1 = \TauOneDefault$ με
$\AblForceExtraPoints{}$ επιπλέον τιμές κοντά στο 1 ($0{,}95$, $0{,}99$, $0{,}999$,
$0{,}9999$), φτάνοντας τα $\AblForceGridPoints{}$ συνολικά σημεία πλέγματος.
Πρόκειται
για ένα σκόπιμα ακραίο, μη ρεαλιστικό καθεστώς λειτουργίας --- δεν προτείνουμε αυτά τα
όρια για ανάπτυξη, τα χρησιμοποιούμε αποκλειστικά για να χαρτογραφήσουμε την ισορροπία
κόστους/ακρίβειας μέχρι το ακραίο της.

Σε ένα ενδιάμεσο σημείο ($\tau_2 = 0{,}99$), το ποσοστό κλήσεων του Επιπέδου~3 ανεβαίνει
στο $\AblForceModerateTierThreePct\%$, με συνολικό \en{Macro F1} $\AblForceModerateFone{}$
και \en{Macro F1} μόνο στο υποσύνολο του Επιπέδου~3 ίσο με
$\AblForceModerateTierThreeFone{}$.
Στο $\tau_2 = 0{,}999$ η καμπύλη πλαταίνει: το
ποσοστό κλήσεων φτάνει $\AblForcePlateauTierThreePct\%$ (το Επίπεδο~2 πέφτει στο
$\AblForcePlateauTierTwoPct\%$), το συνολικό \en{Macro F1} στο $\AblForcePlateauFone{}$,
και το \en{Macro F1} του Επιπέδου~3 στο $\AblForcePlateauTierThreeFone{}$· περαιτέρω
αύξηση σε $\tau_2 = 0{,}9999$ δεν αλλάζει τίποτα --- έχουμε φτάσει σε ένα πλατό, όχι απλά
σε ένα σημείο διαμετρικά κοντά στο 1.

Το Σχήμα~\ref{fig:threshold-ablation-extended} (τρίτο υπο-γράφημα) δείχνει την πλήρη
καμπύλη: το \en{Macro F1} του Επιπέδου~3 ανεβαίνει σταθερά με το ποσοστό κλήσεών του, από
$\TierThreeFone{}$ στην προεπιλεγμένη διαμόρφωση ($\tau_2 = \TauTwoDefault$, μόλις
$\TierThreeCount{}$ ζεύγη) έως $\AblForcePlateauTierThreeFone{}$ στο πλατό, ενώ το
συνολικό \en{Macro F1} της αλυσίδας πέφτει από $\AblDefaultFone{}$ σε
$\AblForcePlateauFone{}$ --- μια πτώση $\AblForceFoneDrop{}$.
Τα δύο ευρήματα δεν
αντιφάσκουν: αφορούν διαφορετικούς πληθυσμούς.
Στα μέτρια $\tau_2$, τα ζεύγη που το
Επίπεδο~2 απορρίπτει είναι απλά ζεύγη στα οποία δεν ήταν πλήρως σίγουρο --- συχνά
ζεύγη που θα επιλύονταν σωστά και εκεί --- οπότε η μεταφορά τους στο ασθενέστερο κατά
μέσο όρο Επίπεδο~3 κοστίζει ακρίβεια.
Καθώς το $\tau_2$ σφίγγει περισσότερο, όμως, τα
ζεύγη που συνεχίζουν να προωθούνται είναι εκείνα όπου το δάσος του Επιπέδου~2 ήταν
λιγότερο ομόφωνο· η σύσταση του πληθυσμού που φτάνει στο Επίπεδο~3 μετατοπίζεται προς
ζεύγη στα οποία είναι συγκριτικά καταλληλότερο, γι' αυτό και το \en{Macro F1} του
υποσυνόλου του ανεβαίνει --- χωρίς αυτό να αναιρεί τη συνολική ζημιά, αφού το Επίπεδο~3
παραμένει κατά μέσο όρο ασθενέστερο του Επιπέδου~2, και κάθε καθαρή μεταφορά πληθυσμού
προς αυτό κοστίζει, όσο καλά και να χειρίζεται πλέον τα συγκεκριμένα ζεύγη.

Ο μηχανισμός πίσω από το πλατό εντοπίζεται στην ίδια την κατανομή εμπιστοσύνης του
Επιπέδου~2.
Από τα \ValSplitSize{} ζεύγη επικύρωσης, $\TierTwoSatEligibleN{}$
($\TierTwoSatEligiblePct\%$) φτάνουν καθόλου στο Επίπεδο~2 --- τα υπόλοιπα επιλύονται ήδη
με εμπιστοσύνη στο Επίπεδο~1.
Από αυτά, το $\TierTwoSatCeilingPct\%$ λαμβάνει \textbf{ομόφωνη}
ψήφο από τον \texttt{\en{PosClassifier}}: όλα τα $\TierTwoSatNEstimators{}$ δέντρα συμφωνούν,
οπότε η εμπιστοσύνη είναι ακριβώς 1{,}0. Η εμπιστοσύνη ενός \en{Random Forest} είναι
κβαντισμένη σε βήματα του $1/\TierTwoSatNEstimators{}$ --- ένα ζεύγος με ομόφωνη ψήφο δεν
μπορεί ποτέ να κλιμακωθεί περαιτέρω, όσο κοντά στο 1 και να τεθεί το $\tau_2$, αφού δεν
υπάρχει τιμή εμπιστοσύνης πάνω από το 1{,}0 να το ξεπεράσει.
Μόλις το
$\TierTwoSatNonSaturatedPct\%$ των ζευγών του Επιπέδου~2 έχει μη-ομόφωνη ψήφο και είναι,
συνεπώς, καν δυνατόν να κλιμακωθεί.
Αυτό εξηγεί επακριβώς το πλατό: το
$\TierTwoSatNonSaturatedPct\%$ επί το $\TierTwoSatEligiblePct\%$ δίνει
$\AblForcePlateauTierThreePct\%$ του συνόλου επικύρωσης --- ακριβώς το ποσοστό κλήσεων
του Επιπέδου~3 στο πλατό.
Το ανώτατο όριο δεν είναι, λοιπόν, ζήτημα επιλογής ορίου· είναι
δομικό χαρακτηριστικό του ίδιου του ταξινομητή του Επιπέδου~2.

Ως σημείο αναφοράς, το ανεξάρτητα εκπαιδευμένο \texttt{\en{EmbeddingClassifier}} (χωρίς
αλυσίδα, εφαρμοζόμενο σε κάθε ζεύγος) έχει \en{Macro F1} $\EmbeddingFone{}$ στο πλήρες
σύνολο επικύρωσης --- το όριο στο οποίο θα τείνει το \en{CascadeLP} αν το $\tau_2$
εξανάγκαζε το 100\% των ζευγών στο Επίπεδο~3.
Το \en{Macro F1} του Επιπέδου~3 στο πλατό
($\AblForcePlateauTierThreeFone{}$) είναι αισθητά υψηλότερο από αυτή την ανεξάρτητη
μέθοδο αναφοράς, αλλά η σύγκριση δεν αποδεικνύει ότι το Επίπεδο~3 «ξεπερνά τον εαυτό του»:
ο πληθυσμός του πλατό είναι ένα \textit{υπόλειμμα} --- τα ζεύγη που τα Επίπεδα~1 και~2
απέρριψαν --- και όχι ένα τυχαίο ή πλήρες δείγμα, οπότε τα δύο νούμερα μετρούν
διαφορετικά πράγματα και δεν συγκρίνονται απευθείας.

Το συμπέρασμα για τη σχεδίαση είναι σαφές: η προεπιλεγμένη τιμή $\tau_2 = \TauTwoDefault{}$
βρίσκεται ήδη κοντά στο βέλτιστο σημείο.
Ο εξαναγκασμός μεγαλύτερου όγκου προς το ακριβό
Επίπεδο~3 δεν προσφέρει κανένα όφελος ακρίβειας και έχει μετρήσιμο κόστος, ενώ το μέγιστο
δυνατό μερίδιο του Επιπέδου~3 είναι δομικά περιορισμένο γύρω στο
$\AblForcePlateauTierThreePct\%$ --- ανεξάρτητα από την επιλογή ορίου.
Αυτό επιβεβαιώνει
όχι μόνο ότι τα φθηνά επίπεδα της αλυσίδας επαρκούν με τις προεπιλεγμένες τιμές, αλλά ότι
επαρκούν σχεδόν πάντα, αφού το μοναδικό σήμα που θα δικαιολογούσε κλιμάκωση --- η
ασυμφωνία μεταξύ των δέντρων του ίδιου του Επιπέδου~2 --- είναι σπάνιο.

\begin{figure}[htbp]
  \centering
  \includegraphics[width=\textwidth]{threshold_ablation}
  \caption{Επέκταση της ανάλυσης ευαισθησίας των ορίων δρομολόγησης. Αριστερά: \en{Macro F1} συναρτήσει του
  $\tau_1$ για το ρεαλιστικό εύρος $\tau_2 \leq 0{,}9$ (\S\ref{sec:ablation}). Κέντρο:
  ποσοστό κλήσεων του Επιπέδου~3 συναρτήσει του $\tau_2$ στο πλήρες επεκταμένο πλέγμα
  ($\tau_1 = \TauOneDefault$), όπου φαίνεται καθαρά το πλατό κοντά στο $\tau_2 = 1$.
  Δεξιά: \en{Macro F1} του υποσυνόλου του Επιπέδου~3 και συνολικό \en{Macro F1} της
  αλυσίδας συναρτήσει του ποσοστού κλήσεων του Επιπέδου~3, στο ίδιο επεκταμένο πλέγμα ---
  η αντίθετη πορεία των δύο καμπυλών αποτυπώνει το εύρημα της παρούσας ενότητας.}
  \label{fig:threshold-ablation-extended}
\end{figure}


\section{Ανάλυση Ζευγών Μηδενικού \en{CN}}
\label{sec:cold-start-results}

Χρησιμοποιώντας τη \texttt{\en{zero\_common\_neighbors\_mask}}
(μηδενικοί κοινοί γείτονες μεταξύ $u$ και $v$) στο σύνολο επικύρωσης, βρίσκουμε
ότι το \ColdStartPct\% των ζευγών (\ColdStartCount{} από \ValSplitSize{},
εξαιρουμένων των αυτοβρόχων) ανήκει στο υποσύνολο μηδενικού \en{CN}.
Μόλις \ColdStartNonCount{} ζεύγη επικύρωσης έχουν έστω έναν κοινό γείτονα, και τα 7 ανήκουν στην
κατηγορία υψηλού \en{CN} και επιλύονται σωστά στο Επίπεδο~1.
Στην πράξη, η
απόδοση ανά επίπεδο στο υποσύνολο μηδενικού \en{CN} ταυτίζεται αριθμητικά με
την συνολική απόδοση ανά επίπεδο του πίνακα~\ref{tab:cascade-tier-results}: το
Επίπεδο~1 χειρίζεται το \TierOnePct\% αυτού του υποσυνόλου με ακρίβεια \TierOneAcc{} και
\en{Macro F1} \TierOneFone{}, το Επίπεδο~2 το \TierTwoPct\% με ακρίβεια \TierTwoAcc{}, και το
Επίπεδο~3 το \TierThreePct\% με ακρίβεια μόλις \TierThreeAcc{}.

Το εύρημα χαρακτηρίζει κυρίως την ακραία αραιότητα του γραφήματος \en{DSAA}: με
μέσο βαθμό \GraphMeanDegree{} (\S\ref{sec:separability-results}), η ύπαρξη κοινού
γείτονα είναι η εξαίρεση.
Δεν αποτελεί μέτρηση του λειτουργικού \en{cold-start}
του \en{CascadeLP}.
Η δρομολόγηση παρακάμπτει το Επίπεδο~1 μόνο όταν τουλάχιστον
ένα άκρο απουσιάζει από το γράφημα, κατάσταση που αναγνωρίζεται από μηδενικές
τιμές και στα τέσσερα δομικά χαρακτηριστικά.
Ένα ζεύγος με μηδενικό \en{CN}
μπορεί να έχει δύο γνωστά άκρα και μη μηδενικό \en{Preferential Attachment},
οπότε εξακολουθεί να αξιολογείται από το Επίπεδο~1.


\section{Ανάλυση του \en{Hard Residual}}
\label{sec:hard-residual}

\subsection{Παραδείγματα Επίλυσης ανά Επίπεδο}
\label{subsec:tier-examples}

Τα συγκεντρωτικά ποσοστά ακρίβειας ανά επίπεδο (\S\ref{sec:cascade-results}) δεν
εξηγούν \textit{γιατί} ένα συγκεκριμένο ζεύγος επιλύεται σε ένα δεδομένο επίπεδο,
ούτε πώς μοιάζει μια λανθασμένη πρόβλεψη σε πραγματικά δεδομένα.
Η ενότητα αυτή παραθέτει ζεύγη πραγματικά αντλημένα από το
\texttt{\en{cascade\_val\_tiers.csv}} (\S\ref{subsec:implementation-architecture}), ένα ανά
συνδυασμό επιπέδου/δυσκολίας/ορθότητας πρόβλεψης, με σκοπό να δείξει ποια πλευρά
των δεδομένων κάθε ταξινομητής βλέπει και πού αποτυγχάνει.
Ο πλήρης κατάλογος
των ζευγών, μαζί με τα ωμά χαρακτηριστικά τους, βρίσκεται στον
Πίνακα~\ref{tab:worked-examples} του Παραρτήματος~\ref{ch:a}.

\subsubsection{Επίπεδο 0 --- Αυτοβρόχοι}

Το Επίπεδο~0 δεν εξετάζει κανένα χαρακτηριστικό· αποφασίζει αποκλειστικά από την
ισότητα $u = v$.
Το ζεύγος \en{(626449, 626449)}, δηλαδή το άρθρο \en{«Young Lake»}
(λίμνη στο πάρκο \en{Adirondack}, Νέα Υόρκη) με τον εαυτό του, έχει ετικέτα θετική
και επιλύεται από τη σύμβαση του Επιπέδου~0, η οποία θεωρεί κάθε αυτοβρόχο θετικό.
Στο σύνολο
εκπαίδευσης όμως υπάρχει ένα μοναδικό αντιπαράδειγμα: το ζεύγος
\en{(74604, 74604)}, το άρθρο \en{«Hartbeespoort Aerial Cableway»}, φέρει ετικέτα
\textit{αρνητική} (\SelfLoopsTrainNeg{} στα \SelfLoopsTrain{} συνολικά, βλ.
\S\ref{sec:leakage-results})\footnote{Δεν έχει βρεθεί εξήγηση γιατί αυτή η
συγκεκριμένη γραμμή φέρει διαφορετική ετικέτα από τις υπόλοιπες 10.220 --- η πιο
πιθανή υπόθεση είναι σφάλμα επισήμανσης κατά την κατασκευή του συνόλου δεδομένων.}.
Ο κανόνας του Επιπέδου~0 προβλέπει θετικά και για τα δύο, άρα αυτή η μία γραμμή
είναι η μοναδική πρόβλεψη σε ολόκληρο το \en{CascadeLP} που είναι λανθασμένη \textit{εκ
κατασκευής} --- όχι λόγω αδυναμίας κάποιου ταξινομητή, αλλά επειδή ο κανόνας
«αυτοσύνδεσμος $\Rightarrow$ θετικό» δεν αφήνει περιθώριο εξαίρεσης για μία ασυνεπή
ετικέτα.

\subsubsection{Επίπεδο 1 --- Δομικοί Ευρετικοί Δείκτες}

\begin{itemize}
  \item \textbf{Ορθή, τετριμμένη λόγω κοινών γειτόνων.} Το γένος εντόμων
  \en{«Molanna»} και το είδος \en{«Molanna ulmerina»} μοιράζονται πολλούς κοινούς
  γείτονες μέσω της κοινής ταξινομικής ιεραρχίας (πρότυπα \en{taxobox}/\en{speciesbox}
  παραπέμπουν και τα δύο σε γένος, οικογένεια, τάξη)· ο αριθμός κοινών γειτόνων
  υπερβαίνει εύκολα το όριο τετριμμένου, το Επίπεδο~1 απαντά ορθά με υψηλή
  αξιοπιστία.
  \item \textbf{Ορθή, υψηλής ομοιότητας κειμένου.} Τα άρθρα \en{«The Golden Mile»}
  (παράκτια περιοχή ερασιτεχνικού ψαρέματος) και \en{«Angling»} (η μέθοδος αλιείας
  καθαυτή) μοιράζονται λεξιλόγιο γύρω από την αλιεία· ο χαρακτηρισμός δυσκολίας
  τα τοποθετεί σε \en{trivial\_high\_textsim}, αλλά και η δομική επικάλυψη αρκεί ώστε
  το Επίπεδο~1 να τα επιλύσει σωστά χωρίς να χρειαστεί κείμενο.
  \item \textbf{Ψευδώς θετικό λόγω προτύπου αποσαφήνισης.} Τα άρθρα \en{«Leaf
  River»} και \en{«Wapello»} είναι και τα δύο σελίδες αποσαφήνισης («may refer
  to:») χωρίς καμία πραγματική σχέση μεταξύ τους.
  Ο χαρακτηρισμός δυσκολίας τα
  βλέπει ως υψηλής ομοιότητας κειμένου (κοινό πρότυπο αποσαφήνισης), αλλά αυτό
  που πραγματικά παραπλανά το Επίπεδο~1 είναι δομικό: σελίδες αποσαφήνισης
  συνδέονται με δεκάδες άσχετους μεταξύ τους κόμβους, πράγμα που τους δίνει
  τεχνητά μεγάλο αριθμό κοινών γειτόνων με \textit{οποιαδήποτε} άλλη σελίδα
  αποσαφήνισης.
  Η πραγματική ετικέτα είναι αρνητική, η πρόβλεψη θετική.
  \item \textbf{Ορθή παρά τον χαρακτηρισμό «δύσκολο».} Η αθλήτρια ενόργανης
  ακροβατικής γυμναστικής \en{«Viktoriya Mikhnovich»} και η χώρα \en{«Belarus»}
  έχουν χαμηλή ομοιότητα κειμένου (βιογραφικό έναντι άρθρου χώρας), οπότε ο
  χαρακτηρισμός δυσκολίας τα βαθμολογεί ως «δύσκολα»· η γειτονιά του γραφήματος όμως
  --- η αθλήτρια συνδέεται με σελίδες της εθνικής ομοσπονδίας και της
  ολυμπιακής αποστολής που με τη σειρά τους συνδέονται με τη χώρα --- δίνει
  επαρκές δομικό σήμα, και το Επίπεδο~1 απαντά σωστά.
  \item \textbf{Ψευδώς θετικό χωρίς προφανή θεματική αιτία.} Το αμφίβιο
  \en{«Plethodontohyla notosticta»} και η φοιτητική αδελφότητα \en{«The Phi Beta
  Kappa Society»} δεν έχουν καμία θεματική σχέση, και η πραγματική ετικέτα είναι
  αρνητική.
  Το Επίπεδο~1 προβλέπει θετικά --- ενδεχομένως λόγω κοινών κόμβων-
  «κόμβων-γέφυρας» (\en{hubs}) στο γράφημα, όπως γενικές κατηγορίες ή έτη ίδρυσης,
  που συνδέουν έμμεσα ζεύγη χωρίς πραγματική σημασιολογική εγγύτητα.
\end{itemize}

\subsubsection{Επίπεδο 2 --- Χαρακτηριστικά Μερών του Λόγου}

\begin{itemize}
  \item \textbf{Ορθή, πρόσωπο--τίτλος.} Ο \en{«Nigel George Paulet, 18th Marquess
  of Winchester»} και ο τίτλος ευγενείας \en{«Marquessate of Winchester»}
  μοιράζονται πυκνή ορολογία γύρω από τη λέξη \en{«Winchester»}/\en{«marquess»}· η
  κατανομή μερών του λόγου (πυκνότητα κύριων ονομάτων, τίτλων) αρκεί για ορθή
  πρόβλεψη.
  \item \textbf{Ψευδώς αρνητικό, αποσαφήνιση προς γεωγραφικό άρθρο.} Η σελίδα
  αποσαφήνισης \en{«Hood River»} θα έπρεπε να συνδέεται με το ομώνυμο γεωγραφικό
  άρθρο (ποταμός στον Καναδά), αλλά η πρόβλεψη είναι αρνητική.
  Η κατανομή μερών
  του λόγου μιας λίστας αποσαφήνισης (κυρίως κύρια ονόματα, ελάχιστα ρήματα)
  διαφέρει αρκετά από αυτήν ενός άρθρου ποταμού ώστε ο ταξινομητής να μην
  αναγνωρίσει τη σχέση αναφοράς --- ένδειξη ότι το επιφανειακό γλωσσολογικό
  σήμα δεν αρκεί εκεί που χρειάζεται η σημασιολογική σχέση «αυτή η σελίδα
  παραπέμπει σε εκείνη».
  \item \textbf{Ορθή, γεωγραφική εγγύτητα.} Το χωριό \en{«No.\ 3 Komarapalayam»}
  (Ταμίλ Ναντού) και ο σιδηροδρομικός σταθμός \en{«Mallur»} χαρακτηρίζονται
  «δύσκολα» λόγω περιορισμένου δομικού σήματος, αλλά μοιράζονται πρότυπα
  υποδομής (\en{infobox settlement}/\en{infobox station}) με παρόμοια κατανομή
  τοπωνυμικών ουσιαστικών, αρκετή ώστε ο \en{PosClassifier} να επιλύσει σωστά τη
  γεωγραφική εγγύτητα.
  \item \textbf{Ψευδώς αρνητικό, ιδίωμα προς ομώνυμο έργο.} Ο αγγλικός
  ιδιωματισμός \en{«Dutch Uncle»} και το ομώνυμο θεατρικό έργο του \en{Simon
  Gray} συνδέονται πραγματικά (το έργο παίρνει τον τίτλο του από τον
  ιδιωματισμό), αλλά η πρόβλεψη είναι αρνητική· η σχέση όρου-προς-έργο-τέχνης
  δεν αποτυπώνεται σε συχνότητες μερών του λόγου, όσο κι αν και οι δύο σελίδες
  μοιράζονται την ίδια φράση-τίτλο.
\end{itemize}

Και στα δύο πρώτα επίπεδα, το επαναλαμβανόμενο μοτίβο είναι ίδιο με αυτό που
περιγράφεται παρακάτω για το Επίπεδο~3: σελίδες αποσαφήνισης και πρότυπα
μορφοποίησης παράγουν ψευδώς θετικά, ενώ σχέσεις αναφοράς/ιδιώματος που
απαιτούν να «διαβαστεί» το νόημα και όχι μόνο η μορφή παράγουν ψευδώς αρνητικά.
Η διαφορά είναι το επίπεδο αφαίρεσης: το Επίπεδο~1 βλέπει μόνο τοπολογία, το
Επίπεδο~2 μόνο κατανομές μερών του λόγου, και κανένα από τα δύο δεν «διαβάζει»
το κείμενο --- αυτό αφήνεται για το Επίπεδο~3.

\subsection[Το Δύσκολο Υπόλοιπο του Επιπέδου 3]{Το Δύσκολο Υπόλοιπο του Επιπέδου~3}
\label{subsec:tier3-hard-residual}

Το υπόλειμμα ζευγών που φτάνει στο Επίπεδο~3 του \en{CascadeLP} \textit{και}
χαρακτηρίζεται «δύσκολο» από τον χαρακτηρισμό διαχωρισιμότητας
(\S\ref{sec:forensics-protocol}) αντιπροσωπεύει το υποσύνολο ζευγών για το οποίο
οι φθηνότερες οικογένειες χαρακτηριστικών γνησίως αποτυγχάνουν: \HardResCount{} ζεύγη στο σύνολο
επικύρωσης (\S\ref{sec:cascade-results}), με ακρίβεια μόλις \HardResAcc{}.

Πριν από την ποιοτική εξέταση, αξίζει να αποκλειστεί μια μη-σημασιολογική εξήγηση: \HardResMissingN{}
από τα \HardResCount{} ζεύγη (\HardResMissingPct\%) περιλαμβάνουν έναν κόμβο με πλήρως ελλιπές κείμενο άρθρου
(τιμή \en{NaN} στο \texttt{\en{nodes.tsv}}), οπότε η ενσωμάτωση \en{sentence-transformer}
δεν διαθέτει ουσιαστικό σήμα εισόδου.
Ένας μόνο κόμβος (\texttt{\en{id=\HardResNodeId{}}}, χωρίς
κείμενο) συμμετέχει σε \HardResNodePairs{} από αυτά τα ζεύγη· σε όλα τα \HardResNodePairs{} η πραγματική ετικέτα είναι
αρνητική, αλλά ο ταξινομητής προβλέπει θετική ετικέτα στο \HardResNodePredPosPct\% από αυτά ---
ουσιαστικά τυχαία πρόβλεψη λόγω απουσίας δεδομένων, όχι αδυναμία σημασιολογικής
συλλογιστικής.
Ωστόσο, η ελλιπής κάλυψη κειμένου \textit{δεν} είναι η κύρια αιτία της
συνολικής κατάρρευσης του Επιπέδου~3: η ακρίβεια στα \HardResCompleteN{} ζεύγη με πλήρες κείμενο και
στις δύο πλευρές είναι \HardResCompleteAcc{}, χειρότερη από την ακρίβεια \HardResIncompleteAcc{} στα \HardResMissingN{} ζεύγη με
ελλιπές κείμενο.
Το γνησίως δύσκολο υποσύνολο είναι εκείνο όπου \textit{υπάρχει}
πλούσιο κείμενο και ο ταξινομητής εξακολουθεί να αποτυγχάνει.

Η χειρωνακτική εξέταση δείγματος από αυτά τα 41 ζεύγη πλήρους κειμένου αναδεικνύει
τρία επαναλαμβανόμενα μοτίβα αποτυχίας:

\begin{itemize}
  \item \textbf{Ομώνυμα σε άσχετους τομείς.} Το άρθρο \en{«FASTRAC 1»} περιγράφει
  διαστημικό δορυφόρο, ενώ το \en{«JCB Fastrac»} περιγράφει γεωργικό ελκυστήρα· τα δύο
  κείμενα μοιράζονται τη λέξη \en{«Fastrac»} αλλά ανήκουν σε εντελώς διαφορετικά
  θεματικά πεδία.
  Η πραγματική ετικέτα είναι θετική (πιθανότατα σύνδεσμος
  αποσαφήνισης λόγω κοινού ονόματος), αλλά η πρόβλεψη είναι αρνητική: η ομοιότητα
  ενσωματώσεων σωστά εντοπίζει τη θεματική απόσταση, λάθος όμως ως προς το αν
  υπάρχει ακμή.
  \item \textbf{Ψευδώς θετικά από κοινό πρότυπο \en{infobox}.} Δύο άρθρα ειδών
  σκαθαριού του γένους \en{Abacetus} (\en{«Abacetus laevigatus»} και ένα δεύτερο
  είδος) μοιράζονται πυκνή ταξινομική ορολογία και πανομοιότυπη δομή προτύπου
  \en{taxobox}, χωρίς καμία πραγματική ακμή μεταξύ τους.
  Η επιφανειακή ομοιότητα
  κειμένου, προερχόμενη από κοινό μορφότυπο και όχι από κοινό περιεχόμενο, οδηγεί σε
  ψευδώς θετική πρόβλεψη.
  \item \textbf{Χαλαρή θεματική συνάφεια.} Το άρθρο αποσαφήνισης \en{«poor me»}
  (τίτλοι τραγουδιών) συνδέεται πραγματικά με το \en{«victim mentality»} (ψυχολογική
  έννοια), μια σχέση νοήματος-προς-τίτλο που απαιτεί βαθύτερη συλλογιστική από την
  ομοιότητα ενσωματώσεων παραγράφου· η πρόβλεψη είναι αρνητική.
\end{itemize}

Στις δύο πρώτες περιπτώσεις η κατεύθυνση του σφάλματος είναι αντίθετη (ψευδώς αρνητικό
έναντι ψευδώς θετικού), αλλά η υποκείμενη αιτία είναι κοινή: η ομοιότητα κειμένου σε
επίπεδο ενσωμάτωσης δεν ταυτίζεται με την ύπαρξη ακμής σε αυτό το σύνολο δεδομένων.
Δύο άρθρα με κοινό όνομα αλλά διαφορετικό αντικείμενο μπορεί να συνδέονται (σύνδεσμος
αποσαφήνισης), ενώ δύο άρθρα με κοινή μορφοποίηση \en{infobox} αλλά ανεξάρτητο
περιεχόμενο μπορεί να μη συνδέονται καθόλου.
Αυτό συνάδει με την τάση υποπρόβλεψης
θετικής ετικέτας που παρατηρήθηκε στη \S\ref{sec:cascade-results} (\TierThreePredPosPct\%
προβλεπόμενα θετικά έναντι \TierThreeTruePosPct\% πραγματικά θετικά): το Επίπεδο~3 τείνει να
ερμηνεύει τη θεματική απόσταση ως απουσία ακμής, κάτι που είναι σωστό στις
περισσότερες περιπτώσεις αλλά όχι στους συνδέσμους αποσαφήνισης/ομωνυμίας.

Το Επίπεδο~3 δεν αποτυγχάνει πάντα σε τέτοια ζεύγη· αξίζει να δούμε και πού
πετυχαίνει, ώστε τα τρία παραπάνω παραδείγματα να μην αφήνουν την εντύπωση
γενικής αδυναμίας.
Το γένος \en{«Abacetus»} εμφανίζεται ξανά, αυτή τη φορά σε
ζεύγος με το εντελώς άσχετο άρθρο \en{«Tidal downsizing»} (υπόθεση σχηματισμού
πλανητών): η πραγματική ετικέτα είναι αρνητική, και εδώ η ομοιότητα
ενσωματώσεων την εντοπίζει σωστά --- σε αντίθεση με το ζεύγος \en{Abacetus}
παραπάνω, δεν υπάρχει κοινό πρότυπο \en{taxobox} που να προκαλεί σύγχυση.
Ομοίως,
το \en{«Niigata Institute of Technology»} (ιδρύθηκε στην πόλη \en{Kashiwazaki}) και
η σελίδα αποσαφήνισης \en{«Kashiwazaki»} συνδέονται πραγματικά, και το Επίπεδο~3
το προβλέπει σωστά χωρίς να χρειάζεται δομική επιβεβαίωση --- η ενσωμάτωση
συλλαμβάνει τη γεωγραφική αναφορά μέσα στο κείμενο του ιδρύματος.

Δύο επιπλέον λανθασμένες προβλέψεις δείχνουν ένα τέταρτο μοτίβο, διαφορετικό
από τα τρία παραπάνω: σχέσεις έργου-προς-δημιουργό ή προσώπου-προς-θεσμικό
κείμενο.
Το μυθιστόρημα \en{«Stormrider»} του συγγραφέα \en{David Gemmell} δεν
συνδέεται, στην πρόβλεψη του Επιπέδου~3, με την ίδια τη σελίδα \en{«David
Gemmell»} --- ενώ η πραγματική ετικέτα είναι θετική.
Η ενσωμάτωση συλλαμβάνει
θεματική εγγύτητα (και τα δύο κείμενα αφορούν φαντασία/λογοτεχνία), αλλά όχι
την ειδικότερη σχέση συγγραφής.
Παρόμοια, ο δικαστής του Ανωτάτου Δικαστηρίου
των \en{ΗΠΑ} \en{«Stephen Breyer»} και ο νόμος \en{«Individuals with Disabilities
Education Act»} έχουν πραγματική θετική ετικέτα, αλλά η πρόβλεψη είναι
αρνητική: η ομοιότητα ενσωμάτωσης δεν αρκεί για να συνδέσει ένα πρόσωπο με ένα
συγκεκριμένο νομοθέτημα, ακόμα κι όταν η σχέση (ο δικαστής έχει εκδώσει
αποφάσεις σχετικές με τον νόμο) είναι πραγματική.
Και στα δύο, το μοτίβο δεν
είναι επιφανειακή ομοιότητα λέξεων όπως στο \en{Fastrac}, αλλά έλλειψη
αναπαράστασης για σχέσεις \textit{δημιουργού--έργου} και
\textit{προσώπου--θεσμικού-κειμένου}, οι οποίες δεν εκφράζονται ρητά στο σώμα
κειμένου του ενός ή του άλλου άρθρου.


\section{Σύγκριση Αποτελεσματικότητας}
\label{sec:efficiency}

Πέρα από την ακρίβεια, η κεντρική υπόσχεση της αρχιτεκτονικής με επίπεδα είναι
η \textit{οικονομία υπολογισμού}: τα φθηνά επίπεδα πρέπει να απορροφούν τη
συντριπτική πλειονότητα των ζευγών, αφήνοντας το ακριβό σημασιολογικό επίπεδο για
ένα ελάχιστο υπόλειμμα.
Η ενότητα αυτή ποσοτικοποιεί αυτή την υπόσχεση.

\subsection{Ρυθμός Επεξεργασίας}
\label{subsec:throughput}

Στην τελική διαμόρφωση μόνο με ευρετικά χαρακτηριστικά (\S\ref{sec:cascade-routing}),
το \en{CascadeLP} επεξεργάζεται τα \ValSplitSize{} ζεύγη του συνόλου επικύρωσης σε
\ThroughputSec~δευτερόλεπτα σε \en{CPU} (ταχύτερη από τρεις εκτελέσεις), ήτοι περίπου
\ThroughputRate{} ζεύγη ανά δευτερόλεπτο ή \ThroughputLatency~ms ανά ζεύγος.
Η μέτρηση αφορά αποκλειστικά
το κόστος δρομολόγησης και ταξινόμησης πάνω σε ήδη υπολογισμένα χαρακτηριστικά· το
κόστος εξαγωγής των χαρακτηριστικών --- κυρίως η κωδικοποίηση \en{sentence-transformer}
του Επιπέδου~3 --- βρίσκεται εκτός αυτής της μέτρησης και αποτελεί το κυρίαρχο
υπολογιστικό κόστος του pipeline.

Στα προεπιλεγμένα όρια, το 99{,}95\% των ζευγών (Επίπεδα~1 και~2,
\TierOneCount{} + \TierTwoCount{} ζεύγη) λαμβάνει απόφαση πριν από το
σημασιολογικό επίπεδο, και μόλις \TierThreeCount{} ζεύγη (\TierThreePct\%) λαμβάνουν
την τελική πρόβλεψη του Επιπέδου~3.
Οι αριθμοί αυτοί μετρούν δρομολόγηση πάνω σε
ήδη υπολογισμένα χαρακτηριστικά.
Υποδεικνύουν πιθανή εξοικονόμηση σε μια μελλοντική
υλοποίηση κατ' απαίτηση, αλλά δεν αποτελούν μέτρηση των κωδικοποιήσεων που
αποφεύχθηκαν στο παρόν πείραμα.

\subsection[Ποσοστό Δρομολόγησης στο Επίπεδο 3 ανά Όριο]{Ποσοστό Δρομολόγησης στο Επίπεδο~3 ανά Όριο}
\label{subsec:tier3-rate}

Το ποσοστό δρομολόγησης στο Επίπεδο~3 (σημασιολογικό επίπεδο βασισμένο σε
\en{sentence-transformer}) καθορίζεται σχεδόν αποκλειστικά από το όριο $\tau_2$: το
$\tau_1$ ρυθμίζει τον καταμερισμό μεταξύ Επιπέδων~1 και~2 και έχει αμελητέα επίδραση
στο υπόλειμμα που κλιμακώνεται πέρα από το Επίπεδο~2.
Ο πίνακας~\ref{tab:tier3-rate} δείχνει το ποσοστό αυτό ως συνάρτηση του $\tau_2$ για $\tau_1 = 0{,}8$ (προεπιλογή).

\begin{table}[htbp]
  \centering
  \begin{tabular}{lrrrrr}
    \toprule
    $\tau_2$ & 0{,}5 & 0{,}6 & 0{,}7 & 0{,}8 & 0{,}9 \\
    \midrule
    Ποσοστό δρομολόγησης στο Επίπεδο~3 & \AblTThreeRateZeroFive\% & \AblTThreeRateZeroSix\% & \AblTThreeRateZeroSeven\% & \AblTThreeRateZeroEight\% & \AblTThreeRateZeroNine\% \\
    \bottomrule
  \end{tabular}
  \caption{Ποσοστό δρομολόγησης στο Επίπεδο~3 ως συνάρτηση του ορίου $\tau_2$
  ($\tau_1 = \TauOneDefault$, σύνολο επικύρωσης n=\ValSplitSize{}).}
  \label{tab:tier3-rate}
\end{table}

Ακόμη και στο πιο επιτρεπτικό σημείο ολόκληρου του πλέγματος ($\tau_1 = 0{,}95$,
$\tau_2 = 0{,}9$), το Επίπεδο~3 χειρίζεται μόλις το \TierThreeCallRateMax\% των ζευγών· στα
προεπιλεγμένα όρια ($\tau_1 = \TauOneDefault$, $\tau_2 = \TauTwoDefault$) το ποσοστό πέφτει στο
\AblTThreeRateZeroSeven\% (\TierThreeCount{} ζεύγη).
Αυτά τα νούμερα αφορούν το ρεαλιστικό εύρος λειτουργίας ($\tau_2 \leq 0{,}9$)· η
\S\ref{subsec:force-tier3} εξετάζει ξεχωριστά τι συμβαίνει αν το $\tau_2$ σφίξει πολύ
πέρα από αυτό, και δείχνει ότι το ποσοστό δρομολόγησης μπορεί μεν να εξαναγκαστεί πολύ
υψηλότερα, αλλά χωρίς κανένα όφελος ακρίβειας.
Το ακριβό σημασιολογικό επίπεδο λειτουργεί, επομένως, ως
σπάνια ενεργοποιούμενη δικλείδα ασφαλείας και όχι ως βασικός φορέας εργασίας ---
συμπέρασμα που ενισχύει την επιλογή να μη γίνεται χαλάρωση του $\tau_2$ πέρα από την
προεπιλογή, δεδομένης της αναξιοπιστίας του επιπέδου στο υπόλειμμα που φτάνει ως εκεί
(\S\ref{sec:hard-residual}).


\section{Ανάλυση Σφαλμάτων}
\label{sec:error-analysis}

Το \en{CascadeLP} κάνει \CascadeTotalErrors{} σφάλματα ταξινόμησης στις \ValSplitSize{} ζεύγη επικύρωσης.
Από
αυτά, το \CascadeErrorsHardPct\% (\CascadeErrorsHardN{} σφάλματα) προέρχεται από ζεύγη που χαρακτηρίζονται \en{hard}
από το πρωτόκολλο χαρακτηρισμού δυσκολίας (\S\ref{sec:forensics-protocol}), και το
\CascadeErrorsTrivialPct\% (\CascadeErrorsTrivialN{} σφάλματα) από τετριμμένα ζεύγη υψηλής ομοιότητας κειμένου --- κανένα
σφάλμα δεν προέρχεται από αυτοβρόχους ή ζεύγη υψηλού \en{CN}, όπου το Επίπεδο~0 και το
Επίπεδο~1 αντίστοιχα είναι αλάνθαστα.

Το μοντέλο αναφοράς \en{TF-IDF + SVM} (\S\ref{sec:svm-baseline}), εκπαιδευμένο και
αξιολογημένο στην ίδια διαίρεση επικύρωσης, κάνει \SvmTotalErrors{} σφάλματα --- περίπου
\CascadeVsSvmRatio{} φορές περισσότερα.
Το προφίλ σφαλμάτων του συγκεντρώνεται στην κατηγορία \en{hard}
σχεδόν στον ίδιο βαθμό με το \en{CascadeLP}: \SvmErrorsHardPct\% από ζεύγη \en{hard} και \SvmErrorsTrivialPct\%
από τετριμμένα ζεύγη υψηλής ομοιότητας κειμένου.
Ο πίνακας~\ref{tab:error-comparison}
συνοψίζει τη σύγκριση.

\begin{table}[htbp]
  \centering
  \begin{tabular}{lrrrr}
    \toprule
    & \multicolumn{2}{c}{\textbf{\en{CascadeLP}}} & \multicolumn{2}{c}{\textbf{\en{SVM}}} \\
    \textbf{Κατηγορία Δυσκολίας} & \textbf{Σφάλματα} & \textbf{\%} & \textbf{Σφάλματα} & \textbf{\%} \\
    \midrule
    \en{hard}                  & \CascadeErrorsHardN{} & \CascadeErrorsHardPct\% & \SvmErrorsHardN{} & \SvmErrorsHardPct\% \\
    Υψηλή Ομοιότητα Κειμένου   & \CascadeErrorsTrivialN{}  & \CascadeErrorsTrivialPct\%  & \SvmErrorsTrivialN{}  & \SvmErrorsTrivialPct\%  \\
    \midrule
    \textbf{Σύνολο}             & \textbf{\CascadeTotalErrors{}} & \textbf{100\%} & \textbf{\SvmTotalErrors{}} & \textbf{100\%} \\
    \bottomrule
  \end{tabular}
  \caption{Σύγκριση προφίλ σφαλμάτων \en{CascadeLP} έναντι \en{SVM} ανά κατηγορία
  δυσκολίας, σύνολο επικύρωσης (n=\ValSplitSize{}).}
  \label{tab:error-comparison}
\end{table}

Το \en{SVM} υποαποδίδει ακόμη και στην κατηγορία «τετριμμένη» υψηλής ομοιότητας
κειμένου: ακρίβεια \SvmHighTextsimAcc{} αλλά \en{Macro F1} μόλις \SvmHighTextsimFone{} σε αυτό το υποσύνολο ---
το ίδιο μοτίβο κατάρρευσης λόγω ανισορροπίας κλάσεων που παρατηρήθηκε στο Επίπεδο~1
του \en{CascadeLP} (\S\ref{sec:cascade-results}).
Στην κατηγορία \en{hard}, το \en{SVM}
πετυχαίνει ακρίβεια \SvmHardAcc{} και \en{Macro F1} \SvmHardFone{}, ενώ το \en{CascadeLP}
πετυχαίνει \CascadeHardAcc{} και \CascadeHardFone{} αντίστοιχα στο ίδιο υποσύνολο --- μια διαφορά που δεν
αποδίδεται σε κοινή χρήση της ομοιότητας \en{TF-IDF}: αυτή αποτελεί τη μοναδική
είσοδο του \en{SVM}, ενώ δεν χρησιμοποιείται από τους ταξινομητές του
\en{CascadeLP}.
Η διαφορά συνδέεται με τον σημαντικά μεγαλύτερο όγκο δεδομένων
εκπαίδευσης που αξιοποιεί το \en{CascadeLP} και με τον συνδυασμό των δομικών
χαρακτηριστικών του Επιπέδου~1, των χαρακτηριστικών \en{POS} του Επιπέδου~2 και
της ομοιότητας \en{Sentence-Transformer} του Επιπέδου~3.
Το Επίπεδο~2 απορροφά το \TierTwoPct\%
των ζευγών επικύρωσης με σχεδόν μηδενικό ποσοστό σφάλματος.
Και τα δύο μοντέλα
συγκεντρώνουν τη μεγάλη πλειονότητα των σφαλμάτων τους στην κατηγορία \en{hard}, σε
σχεδόν πανομοιότυπη αναλογία (\CascadeErrorsHardPct\% έναντι \SvmErrorsHardPct\%) παρά τη 64-πλάσια διαφορά
στο απόλυτο πλήθος σφαλμάτων --- εύρημα που επιβεβαιώνει τη συνέπεια του πρωτοκόλλου
χαρακτηρισμού δυσκολίας ως προγνωστικού δείκτη πραγματικής δυσκολίας ταξινόμησης,
ανεξάρτητα από την οικογένεια μοντέλου.
Ανάλυση ανά σύνολο δεδομένων εκτελείται στην \S\ref{sec:wiki-fresh}, όπου το
\en{CascadeLP} αξιολογείται σε ένα δεύτερο, ανεξάρτητα κατασκευασμένο σύνολο
δεδομένων.




\section{Συγκριτική Ερμηνεία των Δύο Συνόλων}
\label{sec:cross-dataset-discussion}

Τα αποτελέσματα του Πίνακα~\ref{tab:wiki-fresh-results} αλλάζουν την ερμηνεία
της αρχιτεκτονικής χωρίς να αναιρούν τον έλεγχο του \en{DSAA 2023}.
Στο
\en{DSAA}, ένας μικρός αριθμός τιμών του πρώτου άκρου καθορίζει σχεδόν πλήρως
την ετικέτα και ο κανόνας αναζήτησης αναπαράγει τη σχεδόν τέλεια επίδοση.
Το αυτόνομο μοντέλο αναφοράς \en{POS} (0{,}9988) υπερβαίνει μάλιστα οριακά το
\en{CascadeLP} (0{,}9986), άρα η αλυσίδα δεν αποτελεί την καλύτερη μέθοδο στο
αρχικό σύνολο επικύρωσης.
Στο
ανεξάρτητο σύνολο, τα αρνητικά ζεύγη δειγματοληπτούνται από τον ίδιο πληθυσμό
κόμβων και δεν συγκεντρώνονται σε ένα μικρό σύνολο πηγών.
Η επίδοση πέφτει σε
ρεαλιστικότερο επίπεδο, ενώ η σχετική κατάταξη των οικογενειών χαρακτηριστικών
μεταβάλλεται ουσιωδώς.

Η σημαντικότερη διαφορά αφορά τη δομή.
Στο \en{DSAA 2023}, το γράφημα έχει
μέσο βαθμό \GraphMeanDegree{} και τα χαρακτηριστικά γειτονιάς παρέχουν ελάχιστο
σήμα στα περισσότερα ζεύγη.
Στο δεύτερο σύνολο, το γράφημα εκπαίδευσης έχει
μέσο βαθμό \WfMeanDegree{} και ο αυτόνομος δομικός ταξινομητής πετυχαίνει
\en{Macro F1} = \WfStructuralFone{}.
Το \en{CascadeLP} βελτιώνει αυτή την τιμή
σε \WfCascadeFone{} συνδυάζοντας τα δομικά και κειμενικά επίπεδα, ενώ
δρομολογεί το \WfTierOnePct\% των ζευγών στο Επίπεδο~1.
Το εύρημα υποστηρίζει
ότι η αρχιτεκτονική μπορεί να εκμεταλλευτεί πραγματικό δομικό σήμα όταν αυτό
είναι διαθέσιμο.

Η σύγκριση δεν αποδεικνύει ότι το νέο σύνολο αποτελεί καθολικά δυσκολότερο ή
πλήρως αντιπροσωπευτικό πεδίο αξιολόγησης.
Η διάσχιση \en{BFS} (με αφετηρία το
άρθρο \en{«Computer science»}) δημιουργεί ένα
θεματικά και τοπολογικά συνεκτικό υπογράφημα, ενώ τα ομοιόμορφα τυχαία αρνητικά
είναι πιθανό να είναι ευκολότερα από αρνητικά που ταιριάζονται ως προς βαθμό,
θέμα ή κειμενική ομοιότητα.
Επιπλέον, \WfTextEmptyN{} κόμβοι δεν διαθέτουν
κείμενο.
Το διαγνωστικό υποσύνολο των \WfMissingTextCount{} ζευγών επικύρωσης με
τουλάχιστον ένα τέτοιο άκρο δίνει \en{Macro F1} = \WfMissingTextFone{}, αλλά δεν
απομονώνει αιτιακά την επίδραση της απουσίας κειμένου, επειδή τα ίδια ζεύγη μπορεί
να διαθέτουν ισχυρό δομικό σήμα.
Το νέο
πείραμα παρέχει συνεπώς ισχυρότερη ένδειξη για το \en{CascadeLP} από το
\en{DSAA 2023}, αλλά όχι τελική απόδειξη γενίκευσης σε κάθε γράφημα με
κειμενικά χαρακτηριστικά.

Τέλος, τα προηγούμενα αποτελέσματα του διαγωνισμού \en{DSAA 2023} παραμένουν
έγκυρα ως βαθμολογίες στο συγκεκριμένο \en{leaderboard}.
Επειδή όμως οι
δημοσιευμένες μελέτες δεν ελέγχουν τον κανόνα του πρώτου άκρου, οι σχεδόν
τέλειες βαθμολογίες τους δεν αρκούν για να τεκμηριώσουν ότι τα αντίστοιχα
μοντέλα έμαθαν τη σχέση μεταξύ των δύο άρθρων.
Το συμπέρασμα αφορά την
αποδεικτική ισχύ του συνόλου αξιολόγησης και όχι την ορθότητα της υλοποίησης
των προηγούμενων λύσεων.
